{\footnotesize\setlength{\tabcolsep}{4pt}\renewcommand{\arraystretch}{1.0}
\begin{tabular}{@{}>{\raggedright\arraybackslash}p{0.24\linewidth}>{\raggedright\arraybackslash}p{0.72\linewidth}@{}}
\texttt{chunk\_id} & family\_305\_s7 \\
\texttt{provision\_id} & family\_305\_s7 \\
\texttt{act\_id} & family\_305 \\
\texttt{act\_title\_bn} & The Muslim Family Laws Ordinance, 1961 \\
\texttt{act\_year} & 1961 \\
\texttt{provision\_kind} & section \\
\texttt{provision\_no\_bn} & 7 \\
\texttt{provision\_no\_ascii} & 7 \\
\texttt{provision\_title\_bn} & Talaq \\
\texttt{domain} & family \\
\texttt{risk\_tier} & high \\
\texttt{n\_sub} & 1 \\
\texttt{source\_url} & http://bdlaws.minlaw.gov.bd/act-print-305.html \\
\texttt{source\_dataset} & blad \\
\texttt{crawl\_date} & 2025-07-19 \\
\texttt{n\_words} & 226 \\
\texttt{language} & english \\
\texttt{text\_raw} & (1) Any man who wishes to divorce his wife shall, as soon as may be after the pronouncement of talaq in any form whatsoever, give the Chairman notice in writing of his having done so, and shall supply a copy thereof to the wife.(2) Whoever contravenes the provisions of sub-section (1) shall be punishable with simple imprisonment for term which may extend to one year or with fine which may extend to10{[}ten thousand tak … \\
\end{tabular}}
